

\SourceLeaf{472}

\[
(1+\frac{y}{x^{n-1}})(1+\frac{y}{x^{n-3}})(1+\frac{y}{x^{n-5}})\mathinner{.\,.\,.}(1+x^{n-1}y)
\]
Das Product der ersten Hälfte der Factoren
\[
\begin{gathered}=\frac{y^{\frac12n}}{x^{\frac14nn}}.(1+\frac{x^{n-1}}y)(1+\frac{x^{n-3}}y)\\[8pt]\quad{}(1+\frac{x^{n-5}}y)\mathinner{.\,.}(1+\frac{x^3}y)(1+\frac{2}{y}\end{gathered}
\]
wozu noch die übrigen kommen
\[
(1+xy)(1+x^3y(1+x^5y)\mathinner{.\,.\,.}(1+x^{n-1}y)
\]
Da nun
\[
\begin{gathered}A=\frac{1-x^{n+2}}{1-xx}.\frac{1-x^{n+4}}{1-x^4}.\frac{1-x^{n+6}}{1-x^6}\mathinner{.\,.}\frac{1-x^{2n}}{1-x^n}\\[8pt]\quad{}.\frac{y^{\frac12n}}{x^{\frac14nn}}\end{gathered}
\]
wird, so verwandelt sich das erste \textsc{Theorem} in folgendes \textsc{zweite}: \textit{für ein gerades $n$ wird:}
\[
\begin{gathered}
1+\frac{1-x^n}{1-x^{n+2}}.x(y+\frac1y)
\\[8pt]\quad{}+\frac{1-x^n}{1-x^{n+2}}.\frac{1-x^{n-2}}{1-x^{n+4}}.x^4(yy+\frac1{yy})\\[8pt]
\quad\\[8pt]\quad{}+\frac{1-x^n}{1-x^{n+2}}.\frac{1-x^{n-2}}{1-x^{n+4}}.\frac{1-x^{n-4}}{1-x^{n+6}}.x^9(y^3+y^{-3})+\text{etc.}\\[8pt]
=(1+xy)(1+x^3y(1+x^5y)\mathinner{.\,.\,.}(1+x^{n-1}y)
\\[8pt]\quad{}(1+\frac xy)(1+\frac{x^3}y)(1+\frac{x^5}y)\mathinner{.\,.\,.}(1+\frac{x^{n-1}}y)\\[8pt]
\quad\times\frac{1-xx}{1-x^{n+2}}.\frac{1-x^4}{1-x^{n+4}}.\frac{1-x^6}{1-x^{n+6}}\mathinner{.\,.}\frac{1-x^n}{1-x^{2n}}
\end{gathered}
\]
\par\medskip\Needspace{5\baselineskip}\pdfbookmark[1]{Art. 3}{transscendentes-iv-art-3}\centerline{3.}\nobreak\smallskip
Ist $n$ ungerade, so stellt sich die Reihe so dar:
\[
\begin{gathered}
(1+y^n)+\frac{1-a^n}{1-a}t(1+y^{n-2})
\\[8pt]\quad{}+\frac{1-a^n}{1-a}.\frac{1-a^{n-1}}{1-aa}att(1+y^{n-4})+\mathinner{.\,.}\\[8pt]
\quad+\frac{1-a^n}{1-a}.\frac{1-a^{n-1}}{1-aa}\mathinner{.\,.}\frac{1-a^{\frac12(n+3)}}{1-a^{\frac12(n-3)}}\\[8pt]\quad{}a^{\frac18(n-3)(n-5)}t^{\frac12(n-3)}(1+y^3)\\[8pt]
\quad+\frac{1-a^n}{1-a}.\frac{1-a^{n-1}}{1-aa}\mathinner{.\,.}\frac{1-a^{\frac12(n+3)}}{1-a^{\frac12(n-1)}}\\[8pt]\quad{}a^{\frac12.\frac12(n-1)\frac12(n-3)}t^{\frac12(n-1)}(1+y)
\end{gathered}
\]
Machen wir wie vorher $a=xx$ und setzen das Glied, welches hier das letzte ist, $=Bx^{\frac14}.(y^{\frac12}+y^{-\frac12})$, so wird die Reihe
\[
\begin{gathered}
=B\{x^{\frac14}(y^{\frac12}+y^{-\frac12})
\\[8pt]\quad{}+\frac{1-x^{n-1}}{1-x^{n+3}}x^{\frac94}(y^{\frac32}+y^{-\frac32})\\[8pt]
+\frac{1-x^{n-1}}{1-x^{n+3}}.\frac{1-x^{n-3}}{1-x^{n+5}}x^{\frac{25}4}(y^{\frac52}+y^{-\frac52})+\text{etc.}\}
\end{gathered}
\]

\SourceLeaf{473}


welche Reihe aus $\frac12(n+1)$ Gliedern besteht und dann abbricht. Von unserm Product stellen wir die ersten $\frac12(n-1)$ Factoren so dar
\[
\begin{gathered}\frac{y^{\frac12(n-1)}}{x^{\frac14(nn-1)}}.(1+\frac{x^{n-1}}y)(1+\frac{x^{n-3}}y)\\[8pt]\quad{}(1+\frac{x^{n-5}}y)\mathinner{.\,.}(1+\frac{xx}y)\end{gathered}
\]
wozu noch kommt
\[
(1+y)(1+xxy)(1+x^4y)\mathinner{.\,.\,.}(1+x^{n-1}y)
\]
Da nun
\[
\begin{gathered}
B=\frac{1-x^{n+3}}{1-xx}.\frac{1-x^{n+5}}{1-x^4}.\frac{1-x^{n+7}}{1-x^6}\mathinner{.\,.}\frac{1-x^{2n}}{1-x^{n-1}}.
\\[8pt]\quad{}\frac{x^{\frac14nn-n+\frac34}.y^{\frac12(n-1)}y^{\frac12}}{x^{\frac12nn-n+\frac12}.x^{\frac14}}\\[8pt]
=\frac{1-x^{n+3}}{1-xx}\mathinner{.\,.\,.}\frac{1-x^{2n}}{1-x^{n-1}}.\frac{y^{\frac12n}}{x^{\frac14nn}}
\end{gathered}
\]
so ergibt das erste \textsc{Theorem} folgendes \textsc{dritte}: \textit{für ein ungerades $n$ ist}
\[
\begin{gathered}
x^{\frac14}(y^{\frac12}+y^{-\frac12})
\\[8pt]\quad{}+\frac{1-x^{n-1}}{1-x^{n+3}}.x^{\frac94}(y^{\frac32}+y^{-\frac32})
\\[8pt]\quad{}+\frac{1-x^{n-1}}{1-x^{n+3}}.x^{\frac{25}4}(y^{\frac52}+y^{-\frac52})+\text{etc.}\\[8pt]
=x^{\frac14}(y^{\frac12}+y^{-\frac12})(1+xxy)(1+x^4y)(1+x^6y)\mathinner{.\,.\,.}(1+x^{n-1}y)\\[8pt]
\quad\times(1+\frac{xx}y)\quad(1+\frac{x^4}y)\quad(1+\frac{x^6}y)\mathinner{.\,.\,.\,.}(1+\frac{x^{n-1}}y)\\[8pt]
\quad\times\frac{1-xx}{1-x^{n+3}}.\frac{1-x^4}{1-x^{n+5}}.\frac{1-x^6}{1-x^{n+7}}\mathinner{.\,.\,.\,.\,.}\frac{1-x^{n-1}}{1-x^{2n}}
\end{gathered}
\]
\par\medskip\Needspace{5\baselineskip}\pdfbookmark[1]{Art. 4}{transscendentes-iv-art-4}\centerline{4.}\nobreak\smallskip
Wenn man $n$ ins unendliche wachsen lässt, so verwandeln sich die Reihen und Producte des zweiten und dritten Theorems in unendliche Reihen und Producte. In dieser Gestalt ist das \textsc{vierte Theorem}
\[
\begin{gathered}
1+x(y+\frac1y)+x^4(yy+\frac1{yy})+x^9(y^3+\frac1{y^3})+\text{etc. in inf.}\\[8pt]
=(1+xy)(1+\frac xy)(1+x^3y)(1+\frac{x^3}y)\\[8pt]\quad{}(1+x^5y)(1+\frac{x^5}y)\mathinner{.\,.\,.}\\[8pt]
\quad\times(1-xx)(1-x^4)(1-x^6)(1-x^8)\mathinner{.\,.}
\end{gathered}
\]
und \textsc{das fünfte}
\[
\begin{gathered}
x^{\frac14}(y^{\frac12}+y^{-\frac12})+x^{\frac94}(y^{\frac32}+y^{-\frac32})\\[8pt]\quad{}+x^{\frac{25}4}(y^{\frac52}+y^{-\frac52})+\text{etc. in inf.}\\[8pt]
=x^{\frac14}(y^{\frac12}+y^{-\frac12}).(1+xxy)(1+x^4y)(1+x^6y)\mathinner{.\,.}\\[8pt]
\quad\times(1+\frac{xx}y)\quad(1+\frac{x^4}y)\quad(1+\frac{x^6}y)\mathinner{.\,.}\\[8pt]
\quad\times(1-xx)\quad(1-x^4)\quad(1-x^6)\mathinner{.\,.}
\end{gathered}
\]

\SourceLeaf{474}

\par\medskip\Needspace{5\baselineskip}\pdfbookmark[1]{Art. 5}{transscendentes-iv-art-5}\centerline{5.}\nobreak\smallskip
Die Functionen, welche durch das vierte und fünfte Theorem in unendliche Producte entwickelt werden, sind von grosser Wichtigkeit, und es wird gut sein sie hier durch besondere Functionalzeichen zu bezeichnen. Wir schreiben daher
\[
P(x,y)=1+x(y+y^{-1})+x^4(yy+y^{-2})+x^9(y^3+y^{-3})+\text{etc.}
\]
\[
\begin{gathered}R(x,y)=x^{\frac14}(y^{\frac12}+y^{-\frac12})+x^{\frac94}(y^{\frac32}+y^{-\frac32})\\[8pt]\quad{}+x^{\frac{25}4}(y^{\frac52}+y^{-\frac52})+\text{etc.}\end{gathered}
\]
zugleich auch
\[
Q(x,y)=1-x(y+y^{-1})+x^4(yy+y^{-2})-x^9(y^3+y^{-3})+\text{etc.}
\]
wo also $Q(x,y)=P(-x,y)=P(x,-y)$ wird. Der einfachste Werth, welcher $y$ beigelegt werden kann, ist $1$ und da die demselben entsprechenden Werthe unserer Function von besonders grosser Wichtigkeit sind und häufig vorkommen werden, so schreiben wir der Kürze wegen statt $P(x,1),Q(x,1),R(x,1)$ schlechtweg $Px,Qx,Rx$. Wir bemerken noch, dass wo ein Exponent sich blos auf das Argument einer Function bezieht, dieses durch Klammern bezeichnet wird wie $P(x^3)$, ohne Klammern ist immer vorauszusetzen, dass es sich auf die Function bezieht, also $Px^3$ so viel bedeutet wie $(Px)^3$. Also
\[
\begin{gathered}
Px=1+2x+2x^4+\mathord{.}\\[8pt]
Qx=1-2x+2x^4-\mathord{.}\\[8pt]
Rx=2x^{\frac14}+2x^{\frac94}+2x^{\frac{25}4}+\mathord{.}
\end{gathered}
\]
Endlich wollen wir durch das Functionalzeichen $F$ in dieser Abhandlung das unendliche Product ausdrücken
\[
Fx=(1-x)(1-xx)(1-x^3)(1-x^4)\mathinner{.\,.\,.}
\]
In diesen Zeichen erscheinen die beiden letzten Theoreme so
\[
\begin{gathered}\text{4.}\qquad P(x,y)=(1+xy)(1+\frac xy)(1+x^3y)(1+\frac{x^3}y)\\[8pt]\quad{}(1+x^5y)(1+\frac{x^5}y)\mathinner{.\,.\,.} Fxx\end{gathered}
\]
\[
\begin{gathered}\text{5.}\qquad R(x,y)=x^{\frac14}(y^{\frac12}+y^{-\frac12})(1+xxy)(1+\frac{xx}y)\\[8pt]\quad{}(1+x^4y)(1+\frac{x^4}y)\mathinner{.\,.\,.} Fxx\end{gathered}
\]
Und so ist offenbar
\[
\begin{gathered}\text{6.}\qquad Q(x,y)=(1-xy)(1-\frac xy)(1-x^3y)(1-\frac{x^3}y)\\[8pt]\quad{}(1-x^5y)(1-\frac{x^5}y)\mathinner{.\,.\,.} Fxx\end{gathered}
\]





\SourceLeaf{475}

ferner indem man $y=1$ setzt
\[
\begin{gathered}
7.\quad Px=(1+x)^2(1-xx)(1+x^3)^2(1-x^4)\\[8pt]\quad{}(1+x^5)^2(1-x^6)\mathinner{.\,.\,.}\\[8pt]
8.\quad Qx=(1-x)^2(1-xx)(1-x^3)^2(1-x^4)\\[8pt]\quad{}(1-x^5)^2(1-x^6)\mathinner{.\,.\,.}\\[8pt]
9.\quad Rx=2x^{\frac14}(1+xx)^2(1-xx)\\[8pt]\quad{}(1+x^4)^2(1-x^4)(1+x^6)^2(1-x^6)\mathinner{.\,.}
\end{gathered}
\]
Substituirt man hier $1+x=\dfrac{1-xx}{1-x}$, $1+x^3=\dfrac{1-x^6}{1-x^3}$ u.\,s.\,w. so verwandeln diese Ausdrücke sich in folgende
\[
\begin{gathered}
10.\quad Px=\frac{(Fxx)^5}{(Fx)^2(Fx^4)^2}\\[8pt]
11.\quad Qx=\frac{(Fx)^2}{Fxx}\\[8pt]
12.\quad Rx=2x^{\frac14}\frac{(Fx^4)^2}{Fxx}
\end{gathered}
\]
hieraus ergibt sich ferner
\[
\begin{gathered}
13.\quad Px.Qx=(Qxx)^2\\[8pt]
14.\quad Px.Rx=\tfrac12(R\mathord{\surd}{x})^2\qquad\text{oder was dasselbe ist}
\end{gathered}
\]
\[
Pxx.Rxx=\tfrac12(Rx)^2,\qquad Rxx=\frac{(Rx)^2}{2Pxx}
\]
\[
Qxx(Rx)^2=4x^{\frac12}.(Fx^4)^3\qquad\text{also}
\]
\[
\begin{gathered}15.\quad Fx\\[8pt]\quad{}=\mathord{{}^{3}\!\surd}{\frac{Q(x^{\frac12})(Rx^{\frac14})^2}{4x^{\frac18}}}
\\[8pt]\quad{}=\mathord{{}^{3}\!\surd}{\frac{P(x^{\frac12})Q(x^{\frac12})R(x^{\frac12})}{2x^{\frac18}}}
\\[8pt]\quad{}=\mathord{{}^{3}\!\surd}{\frac{(Qx)^2R(x^{\frac12})}{2x^{\frac18}}}\end{gathered}
\]
ferner
\[
\begin{gathered}
16.\quad Px+Qx=2P(x^4)\\[8pt]
17.\quad Px-Qx=2R(x^4)
\end{gathered}
\]
Also durch Multiplication nach 14
\[
18.\quad (Px)^2-(Qx)^2=2(Rxx)^2
\]
Bedeutet ferner $i$ die imaginäre Grösse $\mathord{\surd}{-1}$, so wird
\[
\begin{gathered}
19.\quad Px+iQx=(1+i)Q(ix)\\[8pt]
20.\quad Px-iQx=(1-i)P(ix)
\end{gathered}
\]
Also durch Multiplication

\SourceLeaf{476}

\[
21.\quad (Px)^2+(Qx)^2=2(Pxx)^2
\]
und aus der Multiplication von 18 und 21 mit Zuziehung von 14
\[
22.\quad (Px)^4-(Qx)^4=(Rx)^4
\]
Man sieht also, dass $(Pxx)^2$ das arithmetische Mittel zwischen $(Px)^2$ und $(Qx)^2$ ist, und da nach 13. die Grösse $(Qxx)^2$ das geometrische Mittel zwischen denselben Grössen vorstellt und da (\textit{Theor. attract. el. p.}) wenn zwei Grössenreihen
\[
\begin{array}{llll}
m,&m',&m'',&m'''\mathinner{.\,.\,.}\\[8pt]
n,&n',&n'',&n'''\mathinner{.\,.\,.}
\end{array}
\]
so verbunden sind, dass $m^{(\lambda)}$ immer das arithmetische $n^{(\lambda)}$ das geometrische Mittel zwischen $m^{(\lambda-1)}$ und $n^{(\lambda-1)}$ ist, man die gemeinschaftliche Grenze das arithmetisch geometrische Mittel von $m$, $n$ oder von irgend ein Paar zusammengehörigen Grössen der beiden Reihen nennt, so ergibt sich das höchst wichtige \textsc{Theorem} (23):

\textit{Das Arithmetisch Geometrische Mittel zwischen $(Px)^2$ und $(Qx)^2$ ist allemal $=1$.}

Nach dem, was wir am angezeigten Orte bewiesen haben, ist also auch
\[
24.\quad\text{\textit{das Integral}}\qquad
\int\frac{\mathrm{d}\varphi}{\mathord{\surd}{((Px)^4\cos\varphi^2+(Qx)^4\sin\varphi^2)}}
\]
\textit{von $\varphi=0$ bis $\varphi=2\pi$ ausgedehnt $=2\pi$ oder auch von $\varphi=0$ bis $\varphi=\tfrac12k\pi$, wenn $k$ irgend eine ganze Zahl bedeutet, $=\tfrac12k\pi$.}

\par\medskip\Needspace{5\baselineskip}\pdfbookmark[1]{Art. 6}{transscendentes-iv-art-6}\centerline{6.}\nobreak\smallskip
Um den Zusammenhang des Algorithmus des arithmetisch geometrischen Mittels mit unsern Functionen noch weiter zu entwickeln, bemerken wir zuvörderst, dass wenn
\[
\begin{gathered}\frac{n}{m}=\frac{(Qx)^2}{(Px)^2}\ \text{und}\ m=\mu Px^2,\\[8pt]\quad{}\qquad n=\mu Qx^2,\\[8pt]\quad{}\qquad mm-nn=\mu\mu(R)^4\end{gathered}
\]
gesetzt wird, man hat


\SourceLeaf{477}

\[
\begin{gathered}
m'=\mu(Pxx)^2, n'=\mu(Qxx)^2,\\[6pt]m'm'-n'n'=\mu\mu(Rxx)^4\\[8pt]
m''=\mu(Px^4)^2,n''=\mu(Qx^4)^2,\\[6pt]m''m''-n''n''=\mu\mu(Rx^4)^4\\[8pt]
m'''=\mu(Px^8)^2,n'''=\mu(Qx^8)^2,\\[6pt]m'''m'''-n'''n'''=\mu\mu(Rx^8)^4\\[8pt]
m''''=\mu(Px^{16})^2,n''''=\mu(Qx^{16})^2,\\[6pt]m''''m''''-n''''n''''=\mu\mu(Rx^{16})^4
\end{gathered}
\]
u.\,s.\,w. oder allgemein
\[
\begin{gathered}m^{(\lambda)}=\mu(Px^{2^\lambda})^2,\\[8pt]\quad{}\qquad
n^{(\lambda)}=\mu(Qx^{2^\lambda})^2,\\[8pt]\quad{}\qquad
m^\lambda m^\lambda-n^\lambda n^\lambda=\mu\mu(Rx^{2^\lambda})^4\end{gathered}
\]
und dass offenbar $\mu$ das arithmetisch geometrische Mittel zwischen $m$ und $n$ selbst ist. So bald wir also den Werth von $x$, welcher vorgegebenen Werthen von $m$ und $n$ entspricht, zu bestimmen im Stande sind, werden wir jedes Glied der Reihen
\[
\begin{array}{lll}
m,&m',&m''\mathinner{.\,.}\\[8pt]
n,&n',&n''\mathinner{.\,.}
\end{array}
\]
unmittelbar darstellen, auch die Reihen interpoliren und rückwärts fortsetzen können. Das nächstliegende Mittel $x$ zu bestimmen ist folgendes.

Man sieht leicht, dass die Glieder der Reihe
\[
\begin{gathered}
h=\Bigl(\frac{Rx}{2}\Bigr)^4,\qquad
h'=\Bigl(\frac{Rxx}{2}\Bigr)^2,\qquad
h''=\frac{Rx^4}{2},\qquad
\\[8pt]\quad{}h'''=\Bigl(\frac{Rx^8}{2}\Bigr)^{\frac12},\\[8pt]
h''''=\Bigl(\frac{Rx^{16}}{2}\Bigr)^{\frac14},\qquad
h^{\mathrm{v}}=\Bigl(\frac{Rx^{32}}{2}\Bigr)^{\frac18}\mathinner{.\,.\,.}
\end{gathered}
\]
sich dem $x$ immer mehr nähern, so dass $x$ die Grenze derselben ist. Man hat also
\[
x=h.\frac{h'}{h}.\frac{h''}{h'}.\frac{h'''}{h''}.\frac{h''''}{h'''}\mathinner{.\,.\,.}
\]
Nun ist aber nach 14
\[
\begin{gathered}\frac{h}{h'}=(Pxx)^2\ \text{und so}\ \frac{h'}{h''}=Px^4,\\[8pt]\quad{}\qquad
\frac{h''}{h'''}=(Px^8)^{\frac12},\\[8pt]\quad{}\qquad
\frac{h'''}{h''''}=(Px^{16})^{\frac14}\end{gathered}
\]
Also
\[
\begin{gathered}
x=\frac{h}{(Pxx)^2.Px^4.(Px^2)^{\frac12}.(Px^{16})^{\frac14}\mathinner{.\,.}}\\[8pt]
=\frac{h}{\dfrac{m'}{\mu}.\Bigl(\dfrac{m''}{\mu}\Bigr)^{\frac12}.\Bigl(\dfrac{m'''}{\mu}\Bigr)^{\frac14}.\Bigl(\dfrac{m''''}{\mu}\Bigr)^{\frac18}\mathinner{.\,.\,.}}
\end{gathered}
\]

\SourceLeaf{478}

wir haben folglich
\[
\begin{gathered}25.\quad x=\frac{mm-nn}{16m'\mu\Bigl(\dfrac{m''}{\mu}\Bigr)^{\frac12}.\Bigl(\dfrac{m'''}{\mu}\Bigr)^{\frac14}.\Bigl(\dfrac{m''''}{\mu}\Bigr)^{\frac18}\mathinner{.\,.}}\end{gathered}
\]
wofür man auch schreiben kann
\[
\begin{gathered}x=\frac{mm-nn}{16m'm''}.\Bigl(\frac{m''}{m'''}\Bigr)^{\frac12}.\Bigl(\frac{m'''}{m''''}\Bigr)^{\frac14}.\Bigl(\frac{m''''}{m^{\mathrm{v}}}\Bigr)^{\frac18}\mathinner{.\,.\,.}\end{gathered}
\]
Da die Glieder $m'$, $m''$, $m'''$ sich äusserst schnell der Gleichheit nähern so erhält man $x$ mit grösster Bequemlichkeit.
\begin{center}\rule{0.25\textwidth}{0.3pt}\end{center}





